\documentclass[11pt,a4paper]{article}
\usepackage[utf8]{inputenc}
\usepackage[T1]{fontenc}
\usepackage[portuguese]{babel}
\usepackage{geometry}
\geometry{margin=2.5cm}
\usepackage{listings}
\usepackage{xcolor}
\usepackage{amsmath}
\usepackage{amssymb}
\usepackage{hyperref}
\usepackage{enumitem}
\usepackage{tcolorbox}
\usepackage{tabularx}
\usepackage{booktabs}

\definecolor{codebg}{rgb}{0.95,0.95,0.95}
\definecolor{codegreen}{rgb}{0.1,0.5,0.1}
\definecolor{codepurple}{rgb}{0.5,0.1,0.5}
\definecolor{codegray}{rgb}{0.5,0.5,0.5}

\lstdefinestyle{python}{
    backgroundcolor=\color{codebg},
    commentstyle=\color{codegray}\itshape,
    keywordstyle=\color{codepurple}\bfseries,
    stringstyle=\color{codegreen},
    basicstyle=\ttfamily\small,
    breaklines=true,
    frame=single,
    language=Python,
    showstringspaces=false,
    tabsize=4,
    literate={á}{{\'a}}1 {ã}{{\~a}}1 {é}{{\'e}}1 {í}{{\'i}}1 {ó}{{\'o}}1 {õ}{{\~o}}1 {ú}{{\'u}}1 {ç}{{\c{c}}}1 {Á}{{\'A}}1 {Ã}{{\~A}}1 {É}{{\'E}}1 {Í}{{\'I}}1 {Ó}{{\'O}}1 {Õ}{{\~O}}1 {Ú}{{\'U}}1 {Ç}{{\c{C}}}1 {â}{{\^a}}1 {ê}{{\^e}}1 {ô}{{\^o}}1 {û}{{\^u}}1 {Â}{{\^A}}1 {Ê}{{\^E}}1 {Ô}{{\^O}}1 {Û}{{\^U}}1 {à}{{\`a}}1 {è}{{\`e}}1 {ì}{{\`i}}1 {ò}{{\`o}}1 {ù}{{\`u}}1
}
\lstset{style=python}

\newtcolorbox{exercicio}[1]{
  colback=blue!5!white,
  colframe=blue!40!black,
  title=#1,
  fonttitle=\bfseries
}

\title{\textbf{Data Analysis with Python 2024--2025}\\Parte 3: Matplotlib\\\large Caderno de Exercícios}
\author{Tradução e Adaptação: University of Helsinki --- MOOC.fi}
\date{}

\begin{document}
\maketitle

\section*{Nota Importante}
Este material é uma tradução e adaptação baseada no curso \textit{Data Analysis with Python 2024-2025}, oferecido pela \textbf{The University of Helsinki MOOC Center}.

\tableofcontents
\newpage

\section{Exercícios - Capítulo 3 (Matplotlib)}

\subsection*{Exercício 3.9 -- Vários Gráficos}
\begin{exercicio}{Sumário}
\begin{itemize}
    \item \textbf{Objetivo:} Consulte a documentação de \texttt{plt.plot} para descobrir como especificar explicitamente as coordenadas x. Descubra também como desenhar vários gráficos dentro do mesmo conjunto de eixos.
    \item \textbf{Instruções:} Na função \texttt{main}, faça o gráfico de duas curvas em um único conjunto de eixos. A primeira curva deve utilizar as coordenadas x: \texttt{2, 4, 6, 7} e as coordenadas y: \texttt{4, 3, 5, 1}. A segunda curva deve utilizar as coordenadas x: \texttt{1, 2, 3, 4} e as coordenadas y: \texttt{4, 2, 3, 1}.
    \item \textbf{Detalhes visuais:} Adicione também um título e rótulos para os eixos x e y. No modo não interativo, é necessário chamar \texttt{plt.show()} para exibir a figura.
\end{itemize}
\end{exercicio}

\subsection*{Exercício 3.10 -- Subfiguras}
\begin{exercicio}{Sumário}
\begin{itemize}
    \item \textbf{Objetivo:} Escreva uma função chamada \texttt{subfigures} que crie uma figura contendo duas subfiguras (dois \texttt{axes}, na terminologia do Matplotlib). A função recebe como parâmetro um array bidimensional \texttt{a}.
    \item \textbf{Figura 1:} Na subfigura da esquerda, utilize o método \texttt{plot} para desenhar um gráfico em que as coordenadas x estão na primeira coluna de \texttt{a}, e as coordenadas y estão na segunda coluna de \texttt{a}.
    \item \textbf{Figura 2:} Na subfigura da direita, utilize o método \texttt{scatter} para desenhar os mesmos pontos, mas agora as coordenadas x estão na primeira coluna, as coordenadas y estão na segunda coluna, a cor (\texttt{c}) de cada ponto vem da terceira coluna, e o tamanho (\texttt{s}) de cada ponto vem da quarta coluna. 
    \item \textbf{Instruções Finais:} Para isso, utilize os parâmetros nomeados \texttt{c} e \texttt{s} do método \texttt{scatter}. Teste a função \texttt{subfigures} a partir da função principal.
\end{itemize}
\end{exercicio}

\vspace{2em}
\noindent\textbf{Créditos:} Este conteúdo foi originalmente criado pela \textit{The University of Helsinki MOOC Center} para o curso \textit{Data Analysis with Python}.
\end{document}
